\chapter*{Resumen}

La electroencefalografía (EEG) constituye una herramienta fundamental en neurociencia clínica y experimental, pero sufre de la pérdida frecuente de canales debido a fallas de contacto, artefactos de movimiento o saturación de amplificadores. La interpolación de estos canales faltantes es un paso crítico del preprocesamiento, y tradicionalmente se ha abordado mediante métodos geométricos como los \textit{spherical splines}, implementados en bibliotecas de referencia como MNE-Python. Sin embargo, estos métodos operan exclusivamente en el dominio espacial, ignorando la dinámica temporal de las señales.
%Donde se utiliza conectividad funcional? Si no se utiliza, no hablar de ello.
Esta tesis presenta una evaluación comparativa reproducible y exhaustiva de dieciocho métodos de interpolación de canales EEG, geométricos, estadísticos y basados en Procesamiento de Señales en Grafos (GSP), bajo un protocolo experimental riguroso. El diseño incluye tres bases de datos heterogéneas (PhysioNet de 64 canales, BCI Competition IV 2a de 22 canales y MNE Sample de 60 canales), ciento veinte escenarios de pérdida que combinan modos aleatorio y contiguo con severidades del 10\% al 40\%, y una batería de seis métricas que cubren error de amplitud (MAE, RMSE, SNR), fidelidad morfológica (DTW) y preservación espectral (LSD, Coherencia). Para garantizar una comparación justa, todos los métodos —incluyendo las referencias clásicas— fueron optimizados mediante Optuna con más de quince mil configuraciones paramétricas exploradas, y se incorporó una validación confirmatoria contra la implementación comunitaria de referencia de MNE-Python.

Los resultados muestran una conclusión condicional. En la evaluación confirmatoria balanceada, TRSS fijo mejora la mediana de MAE en 12.4\% frente a MNE-Python, alcanza una tasa de victoria de 72.0\% y mejora NRMSE en 13.9\%. La ventaja aumenta en pérdidas agrupadas o periféricas severas, donde la continuidad temporal compensa la pérdida de información espacial local. No obstante, MNE-Python conserva ventajas importantes: obtiene mejor LSD global, es sustancialmente más rápido y supera a TRSS en señales espacialmente suaves. Se concluye que TRSS es preferible cuando la reconstrucción precisa bajo pérdida no trivial es prioritaria, mientras que MNE sigue siendo la opción robusta para preprocesamiento estándar, análisis espectral sensible a LSD y restricciones de tiempo.

La totalidad del código, las configuraciones experimentales, los registros de optimización y los artefactos de resultados se preservan en el repositorio del proyecto, lo que garantiza la trazabilidad completa de los hallazgos.

\vspace{1em}
\noindent\textbf{Palabras clave:} EEG, canales faltantes, interpolación, Procesamiento de Señales en Grafos (GSP), regularización temporal, TRSS, spherical splines, MNE-Python, optimización bayesiana, Optuna, evaluación comparativa reproducible.